\subsection{pki issue-request: solicitud pública externa}
\label{sec:cli-pki-request}

El titular genera la clave RSA de 3\,072 bits y el CSR PKCS\#10 en su equipo.
La clave del ejemplo queda cifrada; OpenSSL solicita su contraseña localmente.
La solicitud pública se entrega al administrador de la CA:

\begin{lstlisting}[language=bash, numbers=none]
umask 077
openssl genpkey -algorithm RSA -aes-256-cbc \
  -pkeyopt rsa_keygen_bits:3072 -pkeyopt rsa_keygen_pubexp:65537 \
  -out signer-private.pem
openssl req -new -sha256 -key signer-private.pem \
  -out signer.csr.pem \
  -subj '/C=MX/O=Despacho Juridico Demo/CN=Firma documental'
despacho-cli pki issue-request --csr signer.csr.pem \
  --name litigante-demo --json
\end{lstlisting}

El comando admite un CSR PEM de hasta 16~KiB, verifica su prueba de posesión,
el módulo RSA, el exponente 65537 y SHA-256. El identificador de salida admite
de uno a 64 caracteres ASCII seguros. La CA aplica el perfil interno de firma
sin copiar extensiones pedidas y la CLI comprueba que el certificado emitido
conserva la clave pública de la solicitud. País y organización deben coincidir
con la configuración de la CA. Una ruta ya ocupada no se sobrescribe.

La salida JSON contiene \texttt{serial} y \texttt{certificate}; no entrega una
clave privada. El certificado se guarda bajo \texttt{certs/} en el directorio
\texttt{PKI\_CA\_DIR}. El nombre de salida y el nombre común no establecen el
vínculo con una cuenta. Su prueba de posesión, la firma del PDF y el
consentimiento se producen fuera del servidor siguiendo
\rutarepo{docs/individual-signing-guide.md}. El subcomando administrativo de
emisión no inicia una sesión ni firma un documento por el titular.
